\subsubsection{Poisson Equation}

\begin{mystadmonition}{mystGreen}{InteractiveundefinedVisualizationundefined!}Thisundefinednotebookundefinedcontainundefinedinteractiveundefinedvisualizationundefinedthatundefinedneedsundefinedbinderundefinedconnection.undefinedPleaseundefinedclickundefinedtheundefined\textbf{powerundefinedbuttonundefinedonundefinedtheundefinedtopundefinedrightundefinedcorner}undefinedbeforeundefinedyouundefinedstartundefinedreadingundefined!undefinedEstablishingundefinedtheundefinedconnectionundefinedandundefinedbuildingundefinedtheundefinedcontainerundefinedtakesundefinedaundefinedwhile,undefinedsoundefinedclickingundefineditundefinedbeforeundefinedyouundefinedreadundefinedwouldundefinedmakeundefinedsureundefinedtheundefinedinteractiveundefinedcodeundefinedtoundefinedbeundefinedreadyundefinedtoundefinedrunundefinedwhenundefinedyouundefinedscrollundefinedonundefinedtoundefinedthem.

\end{mystadmonition}
\paragraph{2DundefinedPoissonundefinedEquation}

Letundefined$\Omega \subset \mathbb{R}^2$,undefinedweundefinedconsiderundefinedtheundefinedPoissonundefinedequationundefinedwithundefinedhomogenousundefinedboundaryundefinedcondition:

\begin{mystthmdefinition}{2DundefinedPoissonundefinedEquationundefinedwithundefinedHomogenousundefinedBoundaryundefinedCondition}{definition-poissonequation-1}\begin{equation}
\label{2dpoihb}
\begin{aligned}
-\Delta u &= f(x,y) \qquad (x,y)\in \Omega \\
  u(x,y) &= 0 \qquad (x,y)\in \Gamma_1 \\
  \frac{\partial u}{\partial \nu}(x,y) &= 0 \qquad (x,y)\in \Gamma_1
\end{aligned}
\end{equation}

\end{mystthmdefinition}
whereundefined$\Gamma_1\cup \Gamma_2 = \partial \Omega$undefinedsuchundefinedthatundefinedbothundefined$\Gamma_1, \Gamma_2$undefinedisundefinedmeasure-zero.undefined$\nu$undefinedisundefinedtheundefinedoutwardundefinednormalundefinedvector.

\paragraph{SpectralundefinedMethod}

Letundefined
\begin{equation}
u(x,y)=\sum_{-\infty<p,q<\infty} a_{pq}  \exp(i(px +qy))
\end{equation}
Whichundefinedhas:
\begin{equation}
\Delta u = \sum_{-\infty<p,q<\infty} a_{pq}
\end{equation}
Alsoundefinedletundefinedtheundefinedfourierundefinedtransformundefinedofundefined$f(x,y)$undefinedbeundefined$\mathcal F (f(x,y))=F(p,q)$.

\paragraph{Finite-ElementundefinedMethod(FEM)}

\subparagraph{TheundefinedTentundefinedFunction}

Weundefinedconsiderundefinedthe

\subparagraph{ViualizingundefinedtheundefinedTentundefinedBasis}

\begin{mystadmonition}{mystGreen}{Tip}Ifundefinedyouundefinedhaveundefinednotundefineddoundefinedso,undefinedclickundefinedtheundefinedpowerundefinedbuttonundefinedonundefinedtheundefinedtopundefinedrightundefinedcorner.
Afterundefinedeverythingundefinedbuildsundefinedup,undefinedyouundefinedcanundefinedclickundefinedtheundefinedstartundefinedbuttonundefined.undefinedYouundefinedshouldundefinedseeundefinedtwoundefinedslidersundefinedshowingundefinedupundefinedbelow.

\begin{itemize}
\item Theundefinedresolutionundefinedsliderundefineddeterminesundefinedhowundefined``dense''undefinedtheseundefinedtentundefinedbasisundefinedareundefinedinundefinedfillingundefinedtheundefineddomain.undefinedLowerundefinedresolutionundefinedmeansundefinedlargerundefinedandundefinedlessundefinedtentsundefinedandundefinedviceundefinedversa.
\item Theundefinedpositionundefinedsliderundefinedletundefinedyouundefinedadjustundefinedwhichundefinedparticularundefinedbasisundefinedfunctionundefinedyouundefinedareundefinedwatching.undefinedInundefinedotherword,undefinedwhereundefinedtheundefined``tent''undefinedisundefinedlocatedundefinedinundefinedtheundefineddomainundefined$\Omega$
\end{itemize}

\end{mystadmonition}

%  We define:
% :::{prf:definition} Dirchlet Integral
% Let $u\in C^2(\Omega), v\in C^{\infty}(\Omega)$, define the bilinear form on $H^1(\Omega)$ to be
% $$
% a(u,v) := \int_{\Omega}\nabla u \nabla v
% $$
% :::

Weundefinedcanundefinedshowundefinedthat:

\begin{mystthmproposition}{}{proposition-poissonequation-2}Supposeundefined$u \in C^2(\bar \Omega)$undefinedisundefinedaundefinedsolutionundefinedofundefined(\ref{2dpoihb}),undefinedthenundefined$\forall v \in H^1_{\Gamma_1}(\Omega)$,undefinedweundefinedhave:
\begin{equation}
\int_{\Omega}\nabla u \nabla v=\int_{\Omega} fv
\end{equation}

\end{mystthmproposition}
\begin{mystthmproof}{}{proof-poissonequation-3}Weundefinedfirstundefinedshowundefined\textit{Green'sundefinedFormula}(changeundefinedofundefinedvariable):
\begin{equation}
\int_{\Omega}v \Delta u = \int_{\partial \Omega} v \frac{\partial u}{\partial \nu} - \int_{\Omega} \nabla u \nabla v
\end{equation}

\end{mystthmproof}
Thisundefinedmotivatesundefinedusundefinedtoundefineddefineundefinedaundefinedsolutionundefinedinundefinedaundefinedweakerundefinedsense.undefinedSinceundefined$a(u,v)=(f,v)$undefinedisundefinedaundefinednecessaryundefinedconditionundefinedofundefined$u$,undefinedhowundefinedaboutundefinedmakingundefinedeveryundefined$u$undefinedthatundefinedsatisfiesundefinedthisundefinedequationundefinedaundefined``weakundefinedsolution''(orundefinedgeneralizedundefinedsolution)?

Weundefinedfirstundefineddefineundefinedtheundefinedspaceundefinedthatundefinedourundefinedweakundefinedsolutionundefinedwouldundefinedreside:
\begin{equation}
H^1_{\Gamma_1}(\Omega)=\{f\in H^1(\Omega): f=0 \text{ on } \Gamma_1 \}
\end{equation}